\section{Globalne systemy produkcji i uwarunkowania klimatyczne}

Standardy jakości miodu, w~tym stanowisko Apimondia, powstały
na~bazie europejskiego klimatu umiarkowanego i~pszczelarstwa
ramkowego, tymczasem ponad 70\% globalnej produkcji miodu pochodzi
z~regionów klimatycznie, technologicznie i~kulturowo odmiennych
od~Europy~\cite{garcia2025,eudir2001110}. Stosowanie jednolitych
kryteriów morfologicznych i~mikrobiologicznych do~oceny miodów
z~różnych systemów produkcji jest metodologicznie nieadekwatne
i~niesprawiedliwe wobec producentów z~krajów
rozwijających~się~\cite{phiri2022,tadele2025}.

%,----------------------------------------------------------
\subsection{Globalny krajobraz produkcji pszczelarskiej}
%,----------------------------------------------------------

Według danych FAO z~2024~roku globalna liczba zarządzanych
rodzin przekroczyła 101,7~mln, z~czego Azja dominuje
z~44,4\% udziałem, a~Europa posiada 24,7\%~\cite{destatis2026}.
Produkcja miodu w~2022~roku osiągnęła 1,83~mln~ton, przy czym
Azja odpowiadała za~48,2\%, Europa za~22,8\%, obie Ameryki
za~18,5\%, Afryka za~8,5\%~\cite{phiri2022}. Przyjęcie
kryterium zasklepienia jako jedynego wyznacznika dojrzałości
w~praktyce faworyzuje miody europejskie i~dyskryminuje miody
z~systemów tropikalnych i~tradycyjnych~\cite{tadele2025,phiri2022}.

%,----------------------------------------------------------
\subsection{Produkcja afrykańska: wilgotność klimatyczna jako determinanta technologiczna}
%,----------------------------------------------------------

Wilgotność względna powietrza w~strefie tropikalnej Afryki wynosi
typowo 70--90\% RH, w~porze deszczowej przekraczając 95\% RH.
W~takich warunkach pszczoły nie~są biologicznie zdolne do~obniżenia
wilgotności miodu do~$\leq 20\%$ wyłącznie przez naturalną wentylację
ula~\cite{tadele2025}. Pszczelarz staje wobec wyboru: wczesny odbiór
i~dehumidyfikacja technologiczna albo ryzyko fermentacji in situ
lub~plądrowania przez szkodniki. Tadele i~in.~\cite{tadele2025}
klasyfikują dehumidyfikację jako \emph{konieczność technologiczną},
nie~intencję fałszowania, podkreślając, że~alternatywą jest utrata
zbioru lub~produkt sfermentowany. Badania wykazują, że~miód
poddany prawidłowej dehumidyfikacji zachowuje pełny profil
enzymatyczny i~antyoksydacyjny~\cite{lastriyanto2020,spel3upm}.

\subsubsection{Tradycyjne ule kłodowe i koszowe}

W~Etiopii 95,5\% uli stanowią ule tradycyjne, zaledwie 0,2\%
to~ule ramkowe~\cite{ethaau2007}. Paradoksalnie, badanie Gobessa
i~in.~\cite{sciendo2012} wykazało, że~miód z~uli tradycyjnych
charakteryzuje się \emph{istotnie niższą} wilgotnością niż
z~uli nowoczesnych ($p < 0{,}001$), gdyż zbiór odbywa się raz
lub~dwa razy w~roku, co~daje więcej czasu na~naturalne odparowanie.
Jednocześnie ule tradycyjne uniemożliwiają ramkową inspekcję
i~ocenę stopnia zasklepienia bez~niszczenia plastrów, stosowanie
europejskiego kryterium morfologicznego do~tych systemów jest
bezpośrednim przykładem eurocentryzmu normatywnego~\cite{tadele2025}.

%,----------------------------------------------------------
\subsection{Miody bezżądłowych pszczół tropikalnych}
%,----------------------------------------------------------

Pszczoły bezżądłowe (Meliponini, $>500$~gatunków) produkują miód
o~wilgotności typowo 25--35\% i~$a_w = 0{,}60$--$0{,}86$,
odmiennym profilu cukrowym i~wyższej kwasowości niż miody
\textit{Apis mellifera}~\cite{stinglessyeastbrazil2021,stinglesswjava2024}.
Meliponini nie~posiadają mechanizmów odwadniania nektaru
porównywalnych z~pszczołą europejską: mniejsze wole miodowe
i~słabiej rozwinięta wentylacja ula są cechami biologicznymi,
nie~uchybieniami produkcyjnymi. Badanie 17~gatunków
z~Brazylii~\cite{stinglessyeastbrazil2021} wykazało zakres
$a_w = 0{,}60$--$0{,}86$, przy czym drożdże osmofilne izolowane
z~tych miodów charakteryzowały się wysoką osmotolerancją
i~niską zdolnością fermentacyjną, adaptacja biologiczna
do~środowiska o~podwyższonej $a_w$. Miody Meliponini
wymagają odrębnej normalizacji; ocena ich jakości na~podstawie
standardów \textit{Apis mellifera} jest biologicznie
nieuzasadniona~\cite{stinglessyeastbrazil2021,stinglesswjava2024}.

%,----------------------------------------------------------
\subsection{Specyficzne mikroklimaty oraz uwarunkowamna pogodowe strefy umiarkowanej (w tym Europa)} 

Pozyskanie/wirowanie miodu o wilgotności stwarzającej realne ryzyko wystąpienia fermentacji występuje również w strefie umiarkowanej, w tym w Europie. Wiele obszarów z których pozyskiwany jest miód cechuje specyficzny mikroklimat o utrzymującej się podwyższonej wilgotności powietrza. Wpływ na podwyższoną zawartość wody w plastrach mają czynniki pogodowe jak trwające niejednokrotnie przez okres kilkunastu/kilkudziesięciu dni opady.  To właśnie jest powodem, że w stanowisku Apimondii czytamy iż  uzasadnioną praktyką pszczelarską jest pozostawienie miodu  po wirowaniu w otwartych naczyniach celem odparowania nadmiaru wody. Wirowanie miodu o podwyższonej zawartości wody jest więc faktem, podobnie jak działania człowieka zmierzające do usunięcia nadmiaru wody celem stabilizacji mikrobiologicznej. Powszechnie stosowane przez europejskich pszczelarzy jest nie tylko okresowe mieszanie miodu dla zwiększenia szybkości jego parowania, ale wykorzystywane są do tego różne urządzenia techniczne (np. zwiększające powierzchnię rotujące dyski). Rozwój techniczny i automatyzacja procesów widoczne są na wszystkich etapach pozyskiwania miodu i służą nie tylko zwiększeniu wydajności ale i efektywnie podnoszą bezpieczeństwo produktu. Bezsprzecznie, nie mogą być jednak narzędziem fałszowania miodu. Dlatego jasne, dostępne i zweryfikowane kryteria ujęte w akty normalizacyjne winny stać u podstaw oceny takich działań, a w uzasadnionych zakresie winny być precyzyjnie określone praktyki zakazane (np. stosowanie nanofiltracji, żywic jonowymiennych, dehydratacji próżniowej itp).  
%,----------------------------------------------------------
\subsection{Porównanie systemów produkcji}
%,----------------------------------------------------------

Tabela~\ref{tab:hive_systems} zestawia kluczowe cechy trzech
dominujących globalnie typów ula. Tylko ul Langstroth zapewnia
pełną kontrolę nad dojrzewaniem miodu, systemy KTBH
i~ule tradycyjne są w~tym zakresie strukturalnie ograniczone,
co~sprawia, że~miody z~tych systemów mogą nie~spełniać kryterium
zasklepienia przy zachowaniu pełnej dobrej praktyki
pszczelarskiej~\cite{tadele2025,sciendo2012,lrrd2013}.

\begin{table}[H]
\centering
\caption{Porównanie systemów produkcji pszczelarskiej pod~kątem
         kontroli dojrzałości miodu. Źródło: opracowanie własne
         na~podstawie~\cite{tadele2025,sciendo2012,lrrd2013}.}
\label{tab:hive_systems}
\begin{adjustwidth}{-\extralength}{0cm}
\small
\begin{tabularx}{\fulllength}{p{3.5cm}p{3.5cm}p{3.5cm}p{3.5cm}}
\hline
\textbf{Cecha} &
\textbf{Ul Langstroth} &
\textbf{Kenya Top Bar Hive} &
\textbf{Ul tradycyjny} \\
\hline
Kontrola zasklepienia
    & Pełna (inspekcja ramkowa)
    & Częściowa
    & Brak (destrukcyjna) \\
\hline
Dehumidyfikacja in situ
    & Tak
    & Ograniczona
    & Brak \\
\hline
Średni plon [kg/ul/rok]
    & 20--40
    & 15--26~\cite{lrrd2013}
    & 5--15~\cite{tadele2025} \\
\hline
Typowa wilgotność miodu
    & 16--19\%
    & 17--20\%
    & 17--23\% \\
\hline
Ryzyko kolonizacji drożdżami
    & Niskie
    & Umiarkowane
    & Wysokie \\
\hline
Stosowanie globalnie
    & Ameryki, Australia, Europa
    & Afryka Wschodnia
    & Afryka, Azja Pd.-Wsch. \\
\hline
\end{tabularx}
\end{adjustwidth}
\end{table}

%,----------------------------------------------------------
\subsection{Argument normatywny: eurocentryzm standardu morfologicznego}
%,----------------------------------------------------------

Na~podstawie powyższej analizy kryterium pełnego zasklepienia
jako jedynego wyznacznika dojrzałości miodu jest
\emph{eurocentryczne} i~\emph{technicznie nieadekwatne}
z~pięciu powodów: (1)~klimatyczna niemożliwość pełnego
odwodnienia nektaru w~warunkach wilgotności 70--90\%
RH~\cite{tadele2025}; (2)~technologiczna niedostępność oceny
zasklepienia w~ulach tradycyjnych bez~destrukcji
plastrów~\cite{sciendo2012}; (3)~biologiczna odrębność
miodów Meliponini uniemożliwiająca bezpośrednie zastosowanie
standardów \textit{Apis}~\cite{stinglessyeastbrazil2021};
(4)~empiryczna niespójność między stopniem zasklepienia
a~stabilnością mikrobiologiczną miodu~\cite{zhang2021,guo2020}; 
(5)~brak możliwości prowadzenia dowodów, ich archiwizowania i weryfikacji oraz dowodzenia post factum ze stopnia zasklepienia
Konieczne jest przyjęcie zintegrowanego modelu oceny
dojrzałości funkcjonalnej opartego na~mierzalnych
parametrach ($a_w$), nie~na~kryterium
morfologicznym~\cite{tadele2025,codex_stan12,ruegg1981}.